\documentclass[10pt,a4paper]{amsart}
\usepackage[margin=19mm]{geometry}
\usepackage{amsmath,amssymb,mathtools}
\usepackage[scr=rsfs,cal=euler]{mathalfa}
\usepackage{xfrac}
\usepackage[unicode,hidelinks,bookmarks=false]{hyperref}
\newcommand{\ie}{i.e.\ }
\newcommand{\cf}{cf.\ }
\newcommand{\resp}{respectively\ }
\newcommand{\Subsection}{\subsection}
\providecommand{\nobreakdash}{\nobreak\hbox{-}}
\setlength{\emergencystretch}{2em}
\multlinegap0pt
\usepackage{amssymb}
\usepackage{mathtools}
\usepackage{xcolor}
\usepackage{float}
\newtheorem{thm}{Theorem}
\newtheorem{lemma}{Lemma}
\title{Non-Linear Generalization of the DLR Equations: $q$-Specifications and $q$-Equilibrium Measures}
\author{F.H.Haydarov, B.A.Omirov, U.A. Rozikov}
\date{Source: arXiv:2601.06470v2 (CC BY 4.0)}
\begin{document}
\maketitle

\noindent\emph{Source and licence.} Non-Linear Generalization of the DLR Equations: $q$-Specifications and $q$-Equilibrium Measures. Version arXiv:2601.06470v2 (CC BY 4.0), distributed by arXiv under the Creative Commons Attribution 4.0 International licence, http://creativecommons.org/licenses/by/4.0/, as recorded in the arXiv licence metadata for this exact version and shown on the abstract page https://arxiv.org/abs/2601.06470v2. Full licence notice: \texttt{licenses/arxiv-2601.06470-CC-BY-4.0.txt}.\par\smallskip

\noindent\textit{Editorial scope.} Counted unit: the article's thm tt3 (source0.tex lines 2654--2655). The article's complete original proof of it is delivered with it (source0.tex lines 2656--2848, one \texttt{proof} environment of 193 lines). No mathematics is written, repaired or re-derived: the statement and the proof are the article's own lines, in the article's own notation, and no retained line is substituted or re-flowed.  Retained uncounted as the source-local context the proof needs: lines 2650--2650.  Nothing was omitted from any retained span, so the package carries no omission record.  No retained line contains a citation, so the wrapper carries no bibliography and no bibliography entry.  No retained line contains a figure environment, an image-inclusion command or drawing code, so no image is delivered and the package declares no asset dependency.  Editorial formatting only, in addition: an AMS \texttt{amsart} preamble that restates the article's own declarations and macros used by the retained lines, namely the environments \texttt{thm}, \texttt{lemma} on separate counters, together with the standard packages those lines need.  The retained environments print in this document as Theorem 1, Lemma 1 in order of appearance, whereas the article prints the same items with the numbering of its own full document.  The complete native source of the article as distributed in the arXiv e-print for this exact version is bundled unchanged as \texttt{sources/192414/source0.tex}.
\begin{equation}\label{1dh}  H(\sigma_\Lambda; \omega_1, \dots, \omega_q) = -J \sum_{\substack{\{x,y\} \subset \Lambda \\ |x-y|=1}} \sigma_x \sigma_y - h \sum_{x \in \Lambda} \sigma_x - J \sum_{\substack{x \in \Lambda,\, y \in \Lambda^c \\ |x-y|=1}} \sigma_x \, \overline{\omega}_y, \end{equation}

\begin{thm}\label{tt3} For 1D model (\ref{1dh}) if $q\geq 2$, $J>0$, $h=0$ then there is $J_*>0$ such that for each temperature $T \in \left(0, {J\over J_*}\right)$ there are at least three translation-invariant $q$-equilibrium measures.
\end{thm}

\begin{proof} 
{\it The single-site equation}:   Fix the origin $0 \in V$ and take $\Lambda = \{0\}$.
The restriction of the Hamiltonian is
$$
H(\sigma_0; \omega_1, \dots, \omega_q) = -h \sigma_0 - J \sigma_0 \sum_{y \in \partial\{0\}} \overline{\omega}_y,
$$
where $ \partial\{0\} = \{-1, +1\} $. 

Define the boundary function as
$$
\mathbf{b}(\omega_1, \dots, \omega_q) = \sum_{y \in \partial\{0\}} \overline{\omega}_y = \frac{1}{q} \sum_{i=1}^q (\omega_{i,-1} + \omega_{i,1}).
$$

The single-site $ q $-kernel (at inverse temperature $ \beta $) gives 
\begin{equation}\label{ga}
\gamma_{\{0\}}(\sigma_0 = +1 \mid \omega_1, \dots, \omega_q) = \frac{e^{\beta(h + J \mathbf{b})}}{e^{\beta(h + J \mathbf{b})} + e^{-\beta(h + J \mathbf{b})}} = \frac{1 + \tanh(\beta(h + J \mathbf{b}))}{2}.
\end{equation}

{\it Self-consistency equation}:
By compatibility of $\mu$ with the $ q $-specification $ \gamma $, for $ \Lambda = \{0\} $ and the event $ A = \{\sigma : \sigma_0 = +1\} $, we have

$$
\mu(\sigma_0 = +1) = \int_{\Omega^q} \gamma_{\{0\}}(\sigma_0 = +1 \mid \omega_1, \dots, \omega_q) \prod_{i=1}^q \mu(d\omega_i).
$$

Substituting $\gamma_{\{0\}}$ given by (\ref{ga}) we obtain
$$
\mu(\sigma_0 = +1) = \int_{\Omega^q} \frac{1 + \tanh(\beta(h + J \mathbf{b}))}{2} \prod_{i=1}^q \mu(d\omega_i).
$$
Let $ m = \mu(\sigma_0) =\mathbb E_\mu(\sigma_0)=\mu(\sigma_0=1)-\mu(\sigma_0=-1)$. Then $ \mu(\sigma_0 = +1) = \frac{1 + m}{2} $, so
$$
\frac{1 + m}{2} = \frac{1}{2} + \frac{1}{2} \int_{\Omega^q} \tanh(\beta(h + J \mathbf{b})) \prod_{i=1}^q \mu(d\omega_i),
$$
which simplifies to
$$
m = \int_{\Omega^q} \tanh(\beta(h + J \mathbf{b})) \prod_{i=1}^q \mu(d\omega_i). 
$$

This can be written as an expectation over $q$ independent copies of $\mu$:
\begin{equation}\label{Eq}
m = \mathbb{E}_{\mu^{\otimes q}} \left[ \tanh\left( \beta h + \frac{J\beta}{q} \sum_{i=1}^q (\omega_{i,-1} + \omega_{i,1}) \right) \right].
\end{equation}

{\bf Assumption 2.} {\it Let $h=0$, the measure $\mu$ is translation invariant and spins $\omega_{i,\pm 1}$ are independent.} 

Then the measure $ \mu $ is a product measure with spins at different sites independent. 

By $\mathbb{P}$ we denote  the marginal distribution of $\mu^{\otimes q}$ restricted to the specific sites $\omega_{i,-1}$ and $\omega_{i,1}$ for $i = 1, \dots, q$. Thus
$\mathbb{P}$ is a finite-dimensional marginal measure used to compute probabilities related to the sum  $ S = \sum_{i=1}^q (\omega_{i,-1} + \omega_{i,1}) $. 

Since the spins are independent and identically distributed with $ \mathbb{P}(\omega = +1) = \frac{1 + m}{2} $ and $ \mathbb{P}(\omega = -1) = \frac{1 - m}{2} $, we have that $ S $ is the sum of $ 2q $ independent random variables. 

Let $ j $ be the number of $ +1 $ spins among these $ 2q $ spins. Then $ S = 2j - 2q $, and the probability mass function is
$$
\mathbb{P}(S = 2j - 2q) = \binom{2q}{j} \left( \frac{1 + m}{2} \right)^j \left( \frac{1 - m}{2} \right)^{2q - j}.
$$

The expectation (\ref{Eq}) can then be written as
\begin{equation}\label{mf}
m = f(m):=\frac{1}{2^{2q}} \sum_{j=0}^{2q} \tanh\left( \frac{2 J\beta}{q} (j - q) \right) \binom{2q}{j} (1 + m)^j (1 - m)^{2q - j}.
\end{equation}

Equation (\ref{mf}) is a self-consistency condition for the magnetization $ m $. Its solutions correspond to $ q $-equilibrium measures.

{\bf  Case} $ q = 1 $. In this case 
the self-consistency equation is:
$$
m = m \tanh(2J\beta).
$$

For finite $ J\beta $, $ \tanh(2J\beta) < 1 $, so the only solution is $ m = 0 $. Hence, there is exactly one $ q $-equilibrium measure.

{\bf Case} $ q \geq 2 $ and $J>0$. Note that  equation (\ref{mf}) always has the trivial solution $m=0$. We analyze the existence of nontrivial solutions $m\ne 0$ by examining the stability of $m=0$ and determining the critical temperature where a bifurcation occurs.
A bifurcation occurs when $f'(0)=1$.
We have $ f(0) = 0 $, and $ f(-m) = -f(m) $, so $ f $ is odd. It is easy to calculate:
$$
f'(0)=\xi(J\beta):= \frac{2}{2^{2q}} \sum_{j=0}^{2q} (j-q)\,
\tanh\!\left(\frac{2J\beta}{q}(j-q)\right)\binom{2q}{j}$$
\begin{equation}\label{ab}
=2^{2(1-q)} \sum_{j=0}^{q-1} (q-j)\,
\tanh\!\left(\frac{2J\beta}{q}(q-j)\right)\binom{2q}{j}.
\end{equation}
For $J>0$, we consider $x=J\beta>0$, then  $\xi(x)$ is continuous and monotone increasing with maximal (limit as $x\to +\infty$) value 
$$2^{2(1-q)} \sum_{j=0}^{q-1} (q-j)\binom{2q}{j}\geq 2^{1-2q} \left(2^{2q}-\binom{2q}{q}\right)=2\left(1-{1\over 2^{2q}}\binom{2q}{q}\right)\geq 1.25, \ \ q\geq 2.$$ Therefore, denoting unique solution of $\xi(x)=1$ by $J_*>0$, we have  that if $J\beta\in (J_*, +\infty)$ then from (\ref{ab}) we get $|f'(0)|>1$.  For such parameters $J, \beta$ we have three solutions $m=0$, $m=\pm m_*$ of (\ref{mf})  and therefore, there are 3 of $q$-equilibrium measures. 

For  example, for   $ q = 2 $ we have 
\begin{equation}\label{q2}
f'(0) =\xi(J\beta):= \tanh(J\beta) + \frac{1}{2} \tanh(2J\beta).
\end{equation}
It is easy to see that $\xi(x)\in (-1.5, 1.5)$ and 
$J_*\approx 0.647$.

For $q=3$ we have 
$$
f'(0)
= \xi(J\beta):=\frac{3}{16}\Bigg(5\,\tanh\!\Big(\tfrac{2J\beta}{3}\Big)
+ 4\,\tanh\!\Big(\tfrac{4J\beta}{3}\Big)
+ \tanh\!\big(2J\beta\big)\Bigg).
$$
In this case $\xi(x)\in (-1.875, 1.875)$, and 
$J_*\approx 0.5913$.

{\it  Extension to arbitrary finite volumes.}
We now prove that product measure $\mu$ is fixed points of $V_{q,\Lambda}$ for all finite $\Lambda \Subset V$. The key insight is that for any finite volume $\Lambda$, the consistency condition of the $q$-specification and the product structure of $\mu$ ensure that the averaged boundary conditions reduce to constant fields.

Let $\mu$ be one of the product measures identified above with magnetization $m$. For any finite $\Lambda \Subset V$, we need to show
$$
\mu(A) = \int_{\Omega^q} \gamma_\Lambda(A \mid \omega_1, \dots, \omega_q) \prod_{i=1}^q \mu(d\omega_i) \quad \text{for all } A \in \mathcal{F}_\Lambda.
$$

Consider the averaged Hamiltonian. For any configuration $\sigma_\Lambda$, we have
$$
\int_{\Omega^q} H(\sigma_\Lambda; \omega_1, \dots, \omega_q) \prod_{i=1}^q \mu(d\omega_i) = -J \sum_{\substack{\{x,y\} \subset \Lambda \\ |x-y|=1}} \sigma_x \sigma_y - h \sum_{x \in \Lambda} \sigma_x - J \sum_{\substack{x \in \Lambda,\, y \in \Lambda^c \\ |x-y|=1}} \sigma_x \cdot m.
$$

Setting $h = 0$, the averaged kernel becomes
$$
\int_{\Omega^q} \gamma_\Lambda(\sigma_\Lambda \mid \omega_1, \dots, \omega_q) \prod_{i=1}^q \mu(d\omega_i) = \frac{1}{Z_\Lambda(m)} \exp\left( \beta J \sum_{\substack{\{x,y\} \subset \Lambda \\ |x-y|=1}} \sigma_x \sigma_y + \beta J m \sum_{\substack{x \in \Lambda,\, y \in \Lambda^c \\ |x-y|=1}} \sigma_x \right),
$$
where $Z_\Lambda(m)$ is the normalizing constant.

The consistency of the $q$-specification $\{\gamma_\Lambda\}_{\Lambda \Subset V}$ means that for any finite volumes $\Delta \subset \Lambda$, the following compatibility condition holds
$$
\gamma_\Delta(\sigma_\Delta \mid \omega_1, \dots, \omega_q) = \int \gamma_\Lambda(\sigma_\Lambda \mid \omega_1, \dots, \omega_q) \prod_{y \in \Lambda \setminus \Delta} d\sigma_y.
$$

In particular, take $\Delta = \{x\}$. Then the consistency condition becomes
$$
\gamma_{\{x\}}(\sigma_x \mid \omega_1, \dots, \omega_q) = \sum_{\sigma_\Lambda \setminus \{x\}} \gamma_\Lambda(\sigma_\Lambda \mid \omega_1, \dots, \omega_q).
$$

We proof that for any $x \in \Lambda$, the marginal distribution satisfies
\begin{equation}\label{go}
\gamma_\Lambda(\sigma_x \mid \omega_1, \dots, \omega_q) = \gamma_{\{x\}}(\sigma_x \mid \omega_1, \dots, \omega_q).
\end{equation}
Indeed,  take $\Delta = \{x\}$. Then the consistency condition becomes
$$
\gamma_{\{x\}}(\sigma_x \mid \omega_1, \dots, \omega_q) = \sum_{\sigma_\Lambda \setminus \{x\}} \gamma_\Lambda(\sigma_\Lambda \mid \omega_1, \dots, \omega_q).
$$

The right-hand side of this equation is precisely the definition of the marginal distribution of $\sigma_x$ under $\gamma_\Lambda(\cdot \mid \omega_1, \dots, \omega_q)$. That is
$$
\gamma_\Lambda(\sigma_x \mid \omega_1, \dots, \omega_q) := \sum_{\sigma_\Lambda \setminus \{x\}} \gamma_\Lambda(\sigma_\Lambda \mid \omega_1, \dots, \omega_q).
$$ Therefore, we obtain (\ref{go}).

To ensure that this argument applies to our specific model, we verify that the $q$-specification defined by the Hamiltonian (\ref{1dh}) is indeed consistent.

\begin{lemma}
The family $\{\gamma_\Lambda\}_{\Lambda \Subset V}$ defined by
$$
\gamma_\Lambda(\sigma_\Lambda \mid \omega_1, \dots, \omega_q) = \frac{1}{Z_\Lambda(\omega_1, \dots, \omega_q)} \exp\left(-\beta H(\sigma_\Lambda; \omega_1, \dots, \omega_q)\right)
$$
with Hamiltonian given by (\ref{1dh}) forms a consistent $q$-specification.
\end{lemma}

\begin{proof}
Consider finite volumes $\Delta \subset \Lambda$. The Hamiltonian can be decomposed as:
$$
H(\sigma_\Lambda; \omega_1, \dots, \omega_q) = H(\sigma_\Delta; \omega_1, \dots, \omega_q) + H(\sigma_{\Lambda \setminus \Delta}; \omega_1, \dots, \omega_q) + H_{\text{int}}(\sigma_\Delta, \sigma_{\Lambda \setminus \Delta}; \omega_1, \dots, \omega_q),
$$
where $H_{\text{int}}$ contains the interaction terms between $\Delta$ and $\Lambda \setminus \Delta$.

We have
$$
\int \gamma_\Lambda(\sigma_\Lambda \mid \omega_1, \dots, \omega_q) \prod_{y \in \Lambda \setminus \Delta} d\sigma_y = \frac{1}{Z_\Lambda} \sum_{\sigma_{\Lambda \setminus \Delta}} \exp\left(-\beta[H(\sigma_\Delta) + H(\sigma_{\Lambda \setminus \Delta}) + H_{\text{int}}]\right).
$$
This equality can be rewritten as
\begin{equation}\label{DL}
\frac{\exp(-\beta H(\sigma_\Delta))}{Z_\Lambda} \sum_{\sigma_{\Lambda \setminus \Delta}} \exp\left(-\beta[H(\sigma_{\Lambda \setminus \Delta}) + H_{\text{int}}]\right).
\end{equation}

The sum over $\sigma_{\Lambda \setminus \Delta}$ is precisely the partition function for the system on $\Lambda \setminus \Delta$ with boundary conditions determined by $\sigma_\Delta$ and $\omega_1, \dots, \omega_q$. By the definition of  specifications, (\ref{DL}) equals:
$$
\frac{Z_\Delta(\omega_1, \dots, \omega_q)}{Z_\Lambda(\omega_1, \dots, \omega_q)} \gamma_\Delta(\sigma_\Delta \mid \omega_1, \dots, \omega_q),
$$
where the ratio of partition functions ensures proper normalization.

Therefore, we obtain:
$$
\int \gamma_\Lambda(\sigma_\Lambda \mid \omega_1, \dots, \omega_q) \prod_{y \in \Lambda \setminus \Delta} d\sigma_y = \gamma_\Delta(\sigma_\Delta \mid \omega_1, \dots, \omega_q),
$$
which is the desired consistency condition.
\end{proof}

This consistency property is crucial for our proof of Theorem \ref{tt3}. 
It allows us to establish that for any finite $\Lambda \ni x$ and any product 
measure $\mu$ that is a fixed point of $V_{q,\{x\}}$, we have
$$
\int_{\Omega^q} \gamma_\Lambda(\sigma_x \mid \omega_1, \dots, \omega_q) \prod_{i=1}^q \mu(d\omega_i) = \int_{\Omega^q} \gamma_{\{x\}}(\sigma_x \mid \omega_1, \dots, \omega_q) \prod_{i=1}^q \mu(d\omega_i) = \mu(\sigma_x).
$$

This shows that all single-site marginals of $V_{q,\Lambda}(\mu)$ coincide with those of $\mu$. Combined with the product structure of $\mu$ and the 1D nature of the model, this implies that $V_{q,\Lambda}(\mu) = \mu$ for all finite $\Lambda$, completing the proof that $\mu$ is a $q$-equilibrium measure. Theorem \ref{tt3} is proved.
\end{proof}

\end{document}
